\documentclass[12pt]{article}
\usepackage[utf8]{inputenc}
\usepackage[T1]{fontenc}
\usepackage{hyperref}
\title{Tulishi Language Resources for K-12 Curriculum Development}
\author{Compiled by Claude Sonnet 5 (claude-sonnet-5) for Andr\'as Kornai}
\date{}
\begin{document}
\maketitle

\section{Introduction}
Tulishi (ISO 639-3 tey; Glottocode tuli1249) is a Kadu (Kadugli-Krongo) language of South Kordofan, Sudan, with roughly 2,500 to 3,300 speakers. UNESCO classifies it as severely endangered. Overview: \url{https://en.wikipedia.org/wiki/Tulishi_language}.

\section{Primary Educational Resources}
None found. Joshua Project reports that the language has no published written materials: \url{https://joshuaproject.net/people_groups/15626}.

\section{Linguistic and Reference Materials}
Family overview: \url{https://en.wikipedia.org/wiki/Kadu_languages}.
A Leiden University scholarly publication appeared in the results and may concern Kadu languages. Its content was not verified: \url{https://scholarlypublications.universiteitleiden.nl/access/item\%3A2729430/view}.
Omniglot has a Tulishi page: \url{https://omniglot.com/writing/tulishi.htm}.

\section{Religious and Community Resources}
Joshua Project says Bible translation has begun, but no scripture text was found.

\section{Digital and Audio Resources}
None found.

\section{Curriculum Development Context}
Almost nothing exists for classroom use. Any curriculum would have to start from basic documentation and orthography work with the community.
\end{document}